%% kapa: exam=1 round=27 year=2016 period=2 form=A subject=appraisal_law subject_ko=감정평가관계법규 seq=12 number=12
%% source: Q-Net 감정평가사 기출문제 (cst00309) artlSeq=5169402 file=extracted/1008/5169402/제27회 감정평가사 1차 문제/제27회 감정평가사 2교시문제(A형).pdf
%% conversion: kordoc 4.24.51 -> tools/gen1.py (auto); review=proofread-sonnet; flags=none
\begin{question}{12}
국토의 계획 및 이용에 관한 법령상 토지거래계약에 관한 허가신청이 있는 경우의 선매 및 매수청구에 관한 설명으로 옳지 않은 것은?
\begin{choices}
  \item 시장ㆍ군수 또는 구청장은 선매자와 토지소유자 간의 선매협의가 이루어지지 아니한 경우에는 토지수용위원회에 재결을 신청하여야 한다.
  \item 선매자로 지정된 자는 매수가격 등 선매조건을 기재한 서면을 그 지정일부터 15일 이내에 토지소유자에게 통지하여 선매협의를 하여야 한다.
  \item 선매자가 토지를 매수할 때의 가격은 토지거래계약 허가신청서에 적힌 가격이 「부동산가격공시 및 감정평가에 관한 법률」에 따라 감정평가업자가 감정평가한 감정가격보다 높은 경우 감정가격을 기준으로 한다.
  \item 토지거래 불허가처분을 이유로 매수청구를 하는 자는 토지매수청구서를 시장ㆍ군수 또는 구청장에게 제출하여야 한다.
  \item 토지거래 불허가처분을 받은 자는 그 통지를 받은 날부터 1개월 이내에 시장ㆍ군수 또는 구청장에게 해당 토지에 관한 권리의 매수를 청구할 수 있다.
\end{choices}
\end{question}
\answer{1}
